\documentclass[11pt,a4paper]{article}
\usepackage[T1]{fontenc}
\usepackage[utf8]{inputenc}
\usepackage{libertinus}
\usepackage{microtype}
\usepackage[margin=25mm,headheight=15pt]{geometry}
\usepackage{amsmath,amssymb,amsthm,mathtools}
\usepackage{booktabs,array,tabularx}
\usepackage{enumitem}
\usepackage{xcolor,xurl}
\usepackage{hyperref}
\usepackage{fancyhdr}
\definecolor{ink}{HTML}{183249}
\definecolor{accent}{HTML}{176B70}
\hypersetup{colorlinks=true,linkcolor=ink,citecolor=accent,urlcolor=accent,
 pdftitle={Gaussian Confounding and Sharp Recovery of Uniform Convolution Factors},
 pdfauthor={Research report prepared for Vladimir Reshetnikov},
 pdfsubject={Shifted moment inversion, uniform convolutions, sharp minimax rates}}
\pagestyle{fancy}\fancyhf{}
\fancyhead[L]{\small Gaussian confounding and uniform factors}
\fancyhead[R]{\small\thepage}
\renewcommand{\headrulewidth}{0.3pt}
\setlength{\parskip}{0.3em}
\setlength{\emergencystretch}{2em}
\setlist{itemsep=3pt,topsep=4pt}
\allowdisplaybreaks[2]
\numberwithin{equation}{section}
\newtheorem{theorem}{Theorem}[section]
\newtheorem{lemma}[theorem]{Lemma}
\newtheorem{proposition}[theorem]{Proposition}
\newtheorem{corollary}[theorem]{Corollary}
\theoremstyle{definition}
\newtheorem{definition}[theorem]{Definition}
\newtheorem{example}[theorem]{Example}
\newtheorem{problem}{Research question}
\theoremstyle{remark}
\newtheorem{remark}[theorem]{Remark}
\newcommand{\R}{\mathbb R}
\newcommand{\E}{\mathbb E}
\newcommand{\Pp}{\mathbb P}
\newcommand{\Law}{\mathcal L}
\newcommand{\TV}{d_{\mathrm{TV}}}
\newcommand{\dd}{\,\mathrm d}
\newcommand{\norm}[1]{\left\lVert#1\right\rVert}
\newcommand{\abs}[1]{\left|#1\right|}
\newcommand{\sinc}{\operatorname{sinc}}
\newcommand{\clip}{\operatorname{clip}}
\newcommand{\kap}{\kappa}
\newcommand{\As}{\mathcal A_m}
\newcommand{\Xs}{\mathcal X_m}
\newcommand{\risk}{\mathcal R}
\newcommand{\He}{\operatorname{He}}
\newcommand{\supp}{\operatorname{supp}}
\title{\textbf{Gaussian Confounding and Sharp Recovery\\of Uniform Convolution Factors}\\[0.55em]
\large Shifted power sums, minimax estimation, and detection\\above a Fabius--Rvachev background}
\author{Research report prepared for Vladimir Reshetnikov}
\date{29 September 2026}
\begin{document}
\maketitle
\begin{abstract}
Let a known bounded random variable be convolved with at most $m$ centered
uniform factors and a Gaussian whose variance may be unknown. We determine
the exact global minimax rates for recovering the ordered uniform half-lengths
from independent observations. The rate is $n^{-1/(4m)}$ when the Gaussian
variance is known and $n^{-1/(4m+4)}$ when it is unknown. In the latter model,
the optimal rate for the Gaussian variance itself is $n^{-1/(2m+2)}$.
The constants may depend on the fixed capacity $m$, the background, and a
compact positive variance interval.

The algebraic engine is a sharp shifted-power-sum inequality: the moments of
orders $r+1,\ldots,r+m$ determine $m$ ordered nonnegative nodes with optimal
inverse H\"older exponent $1/(m+r)$. We prove this directly using sign changes
of an integer-valued counting function, and construct exact moment-preserving
curves attaining the exponent. Gaussian compensation converts the missing
first power sum into a genuine loss of statistical information. A weighted
likelihood expansion supplies matching Hellinger and chi-square asymptotics,
closing the gap between algebraic moment cancellation and statistical lower
bounds.

Detection has a different scale: the presence of any nonzero uniform factor
has critical separation $n^{-1/8}$ with unknown Gaussian variance, independent
of $m$, versus $n^{-1/4}$ with known variance. The Rvachev up law is an explicit
admissible background. The article includes proofs, a feasible moment estimator,
exact finite algebraic checks, a source comparison, and a research agenda.
These are conventional mathematical proofs, not a Lean verification or a claim
of worldwide publication priority.
\end{abstract}
\newpage
\setcounter{tocdepth}{1}
\tableofcontents
\newpage

\section{Research target and relation to ProveIt}
\label{sec:scope}

\subsection{The extension addressed here}
A convolution factorization can be identifiable while its recovery from noisy
data is extremely unstable. The ProveIt repository develops this question for
Fabius--Rvachev laws and more general uniform convolutions. Its recent report
\emph{Sharp Stability Strata for Finite Fabius--Rvachev Deconvolution}
\cite{RepoFinite} studies at most $m$ unknown uniform factors over a known smooth
compactly supported background. It gives deterministic total-variation inverse
exponents, explicitly distinguishes pairwise and fixed-base stability, and
explicitly does not establish a statistical minimax lower bound. The repository
index records that report as an unreviewed research intake rather than a
Lean-verified theorem \cite{RepoIndex}.

We change the statistical experiment in an essential way: there is a
\emph{positive Gaussian convolution}, and its variance can be an unknown
nuisance parameter. Gaussian noise contributes only to the second cumulant.
Consequently, it can conceal the variance of small uniform factors while
leaving their higher cumulants visible. How much statistical information does
this remove? Can one recover the noise variance and all the factors at matching
upper and lower rates? Is detecting a factor as hard as recovering the entire
factorization? These are the concrete questions answered below.

Our conclusions are not asserted for the unsmoothed experiment of the preceding
report. Compact-support boundaries create additional likelihood issues there.
Positive Gaussian smoothing is a hypothesis, not a disposable technical
convenience. Conversely, we require only that the known background be bounded;
it need not have a density or be symmetric.

\subsection{Classical ingredients and the contribution boundary}
Cumulants, Newton identities, Vandermonde matrices, and the geometry of singular
moment inversion are classical. Prony-map singularities provide closely related
algebraic context \cite{BatenkovYomdin}. Moment fitting and moment-matching
statistical lower bounds have an extensive literature for finite mixtures;
Heinrich--Kahn and Wu--Yang are relevant examples \cite{HeinrichKahn,WuYang}.
Those mixture models have unknown component locations and usually unknown
mixing weights. Our model instead has a \emph{sum of independent uniforms},
with fixed unit multiplicities in its power sums. Mixture theorems are not
being applied as though they directly proved the rates below.

The results developed here are a sharp consecutive shifted-moment inequality,
an explicit variance-confounding moment flow, a weighted likelihood expansion,
and matching estimation and testing bounds for this convolution model. All
proofs needed for those conclusions are supplied. The source comparison is
focused, not an exhaustive worldwide novelty audit; no claim of publication
priority is made. None of the new proofs has been independently refereed or
checked in Lean or Rocq.

The repository snapshot used for the principal comparison is
\begin{center}
\small\texttt{c130dba623c420551d90db41dae7b3d9cff72dfd}.
\end{center}
The accompanying \texttt{SOURCE\_AUDIT.md} records the exact paths and scope.
Search-index hits at other commits were used for discovery, not silently
identified with this snapshot. No repository files were modified.

\section{Model and main statistical results}
\label{sec:model}

Fix an integer $m\ge1$ and a known random variable $Z$ with $|Z|\le L$ almost
surely. Let $U_1,\ldots,U_m$ be independent uniforms on $[-1,1]$, and let $G$ be
a standard normal, all independent of $Z$. Write
\begin{equation}
 \As=\{a\in[0,1]^m:a_1\le\cdots\le a_m\},\qquad
 X_{a,v}=Z+\sum_{i=1}^m a_iU_i+\sqrt v\,G,
 \quad P_{a,v}=\Law(X_{a,v}).
 \label{eq:model}
\end{equation}
Zero entries pad the list to capacity $m$; they do not represent observable
factors. The variance parameter is either fixed at a known $v_0>0$, or ranges
in a fixed interval
\begin{equation}
 0<\underline v<\overline v<\infty.
 \label{eq:varianceinterval}
\end{equation}
We observe $n$ independent copies of $X_{a,v}$. Ordered maximum error is the
loss for $a$, and absolute error is the loss for $v$. Thus, in the unknown
variance model, define
\begin{align}
 \risk^a_n&=\inf_{\widehat a}\sup_{a\in\As,\ v\in[\underline v,\overline v]}
       \E_{a,v}\norm{\widehat a-a}_\infty,\\
 \risk^v_n&=\inf_{\widehat v}\sup_{a\in\As,\ v\in[\underline v,\overline v]}
       \E_{a,v}|\widehat v-v|.
 \label{eq:risks}
\end{align}
Estimators may be measurable functions of the entire sample. Let
$\risk^{a,\mathrm{known}}_n$ denote the analogous scale risk with $v=v_0$ fixed.
All asymptotic constants below can depend on fixed $m,L,Z$ and the variance
bounds (or $v_0$), but not on $n$ or the varying parameters.

\begin{theorem}[Sharp global estimation rates]
\label{thm:main}
There are positive constants $c,C$ and $n_0$ such that, for $n\ge n_0$,
\begin{align}
 c n^{-1/(4m+4)}\le\risk^a_n&\le C n^{-1/(4m+4)},
 \label{eq:unknownrate}\\
 c n^{-1/(2m+2)}\le\risk^v_n&\le C n^{-1/(2m+2)},
 \label{eq:variancerate}\\
 c n^{-1/(4m)}\le\risk^{a,\mathrm{known}}_n&\le C n^{-1/(4m)}.
 \label{eq:knownrate}
\end{align}
A single moment-fitting procedure attains both upper bounds in the unknown
variance model.
\end{theorem}

\begin{center}
\begin{tabular}{@{}llll@{}}
\toprule
Capacity & Known-variance scales & Unknown-variance scales & Unknown variance\\
\midrule
$m=1$ & $n^{-1/4}$ & $n^{-1/8}$ & $n^{-1/4}$\\
$m=2$ & $n^{-1/8}$ & $n^{-1/12}$ & $n^{-1/6}$\\
$m=3$ & $n^{-1/12}$ & $n^{-1/16}$ & $n^{-1/8}$\\
General $m$ & $n^{-1/(4m)}$ & $n^{-1/(4m+4)}$ & $n^{-1/(2m+2)}$\\
\bottomrule
\end{tabular}
\end{center}

For detection, compare the null $a=0$ with alternatives
$\norm a_\infty\ge h$. The variance is allowed to vary under both hypotheses
when it is unknown. Write $\mathcal T_n(h)$ for the infimum, over tests, of
the sum of the maximal type-I and type-II errors.

\begin{theorem}[Detection boundary]
\label{thm:detectionintro}
In the unknown variance model, there is a test with
$\mathcal T_n(h)\le C/(nh^8)$. If $h_n\to0$ and $nh_n^8\to0$, then
$\mathcal T_n(h_n)\ge1-o(1)$. With known variance, the corresponding statements
hold with $h^4$ in place of $h^8$.
\end{theorem}

Thus consistent detection is possible when $nh^8\to\infty$ and impossible
when $nh^8\to0$ in the unknown-variance experiment. The critical order is
$n^{-1/8}$, independently of $m$. For $m>1$, recovering every half-length to
its worst-case optimal precision is strictly harder than detecting the
presence of at least one factor. Proofs appear in Sections~\ref{sec:upper},
\ref{sec:lower}, and~\ref{sec:testing}.

\section{Cumulants and the missing first power sum}
\label{sec:cumulants}

For a vector $x=(x_1,\ldots,x_m)$, let $p_k(x)=\sum_i x_i^k$ and put
$x_i=a_i^2$. All random variables in our model have moment generating
functions in a neighborhood of the origin. Cumulants are defined by
\begin{equation}
 \log\E e^{tX}=\sum_{j\ge1}\kap_j(X)\frac{t^j}{j!}.
\end{equation}
For $U$ uniform on $[-1,1]$, $\E e^{tU}=\sinh(t)/t$. Its even cumulants are
\begin{equation}
 c_k:=\kap_{2k}(U)=\frac{2^{2k}B_{2k}}{2k}
 =(-1)^{k+1}\frac{(2k)!\zeta(2k)}{k\pi^{2k}},
 \qquad k\ge1,
 \label{eq:ck}
\end{equation}
where the Bernoulli convention is $B_2=1/6$. The formula follows by integrating
the Taylor series of $\coth t-1/t$, or by expanding the product for
$\sinh(t)/t$. In particular,
\begin{equation}
 c_1=\frac13,\qquad c_2=-\frac{2}{15},\qquad c_3=\frac{16}{63};
 \quad c_k\ne0\text{ for every }k.
\end{equation}
Independence and scaling give the exact identities
\begin{align}
 \kap_2(X_{a,v})-\kap_2(Z)&=v+\frac13p_1(x),
 \label{eq:varianceidentity}\\
 \kap_{2k}(X_{a,v})-\kap_{2k}(Z)&=c_kp_k(x),\qquad k\ge2.
 \label{eq:cumulantidentity}
\end{align}
All odd cumulants of the added noise are zero. If $v$ is known, $p_1$ is
observed through the second cumulant. If $v$ is unknown, only the combination
$v+p_1/3$ is observed there. The available uncontaminated consecutive power
sums begin at $p_2$, not $p_1$.

\subsection{The Fabius--Rvachev specialization}
Arias de Reyna discusses the smooth compactly supported up function and its
infinite product representation \cite{Arias}. For our normalization it suffices
to define the background probabilistically:
\begin{equation}
 Z_{\mathrm{up}}=\sum_{j=1}^{\infty}2^{-j}U_j.
 \label{eq:up}
\end{equation}
The series converges absolutely for every outcome and has support in $[-1,1]$.
Its characteristic function is
\begin{equation}
 \E e^{itZ_{\mathrm{up}}}=\prod_{j=1}^{\infty}\sinc(2^{-j}t),
 \qquad \sinc t=\frac{\sin t}{t},\quad\sinc0=1.
\end{equation}
Bounded convergence gives the product, and locally convergent logarithms give
\begin{equation}
 \kap_{2k}(Z_{\mathrm{up}})=\frac{c_k}{2^{2k}-1}.
 \label{eq:upcumulants}
\end{equation}
In particular, its variance is $1/9$ and its fourth cumulant is $-2/225$.
These formulas make the required background cumulants explicit. Our proofs
use boundedness, not any unproved regularity claim about the up density.

\section{A sharp inverse theorem for shifted power sums}
\label{sec:shifted}

Let $\Xs=\{x\in[0,1]^m:x_1\le\cdots\le x_m\}$. For an integer $r\ge0$ define
\begin{equation}
 E_r(x,y)=\max_{r+1\le k\le r+m}|p_k(x)-p_k(y)|.
\end{equation}

\begin{theorem}[Consecutive shifted power sums]
\label{thm:shifted}
For all $x,y\in\Xs$ and all integers $r\ge0$,
\begin{equation}
 \norm{x-y}_\infty\le 8 E_r(x,y)^{1/(m+r)}.
 \label{eq:shifted}
\end{equation}
The exponent $1/(m+r)$ cannot be replaced by any larger exponent in a uniform
bound on $\Xs$. Consequently, the map
$x\mapsto(p_{r+1}(x),\ldots,p_{r+m}(x))$ is injective.
\end{theorem}

\begin{proof}[Proof of the upper bound]
Set
\[
 C(t)=\#\{i:x_i\le t\}-\#\{i:y_i\le t\},\qquad 0<t<1.
\]
This is an integer-valued step function. Its nonzero sign changes at most
$m-1$ times. To see this even when atoms coincide, split its jumps into unit
up-steps and down-steps. There are $m$ of each, and the resulting walk starts
and ends at zero. Each excursion from zero consumes at least one up-step and
one down-step, so there are at most $m$ excursions. Consecutive nonzero sign
blocks of $C$ can only decrease this count. Intermediate steps at a common
atom occupy no interval and do not create an extra sign block of $C$.

Choose points $z_1,\ldots,z_s\in[0,1]$ between consecutive nonzero sign blocks
of opposite sign, with $s\le m-1$. There is a sign choice such that
\[
 Q(t)=\pm\prod_{j=1}^s(t-z_j),\qquad Q(t)C(t)\ge0
\]
almost everywhere. If there are no sign changes, $Q$ is the appropriate
constant $1$ or $-1$. The sum of the absolute values of the coefficients of
$Q$ is at most $2^s$.

Put $h=\norm{x-y}_\infty$. There is nothing to prove for $h=0$. By interchanging
$x,y$ if necessary, some $i$ satisfies $y_i-x_i=h$. For $x_i<t<y_i$,
$C(t)\ge1$. No $z_j$ lies in this open interval. On its middle half, which
has length $h/2$, we therefore have $t\ge h/4$ and $|t-z_j|\ge h/4$ for every
$j$. It follows that
\begin{equation}
 \int_0^1 t^r Q(t)C(t)\dd t
 \ge \frac h2\left(\frac h4\right)^{r+s}.
 \label{eq:lowercount}
\end{equation}
For $k\ge1$, direct integration of the counting functions gives
\begin{equation}
 p_k(x)-p_k(y)=-k\int_0^1t^{k-1}C(t)\dd t.
 \label{eq:countmoment}
\end{equation}
Writing $Q$ in monomials and using $s\le m-1$ yields
\begin{equation}
 \abs{\int_0^1t^rQ(t)C(t)\dd t}
 \le\frac{2^s}{r+1}E_r(x,y).
\end{equation}
Thus $h^{r+s+1}\le 2^{2r+3s+1}E_r/(r+1)$. Since $0<h\le1$ and
$s\le m-1$, we obtain
\[
 h^{m+r}\le\frac{2^{2r+3m-2}}{r+1}E_r(x,y)
 \le8^{m+r}E_r(x,y),
\]
which proves~\eqref{eq:shifted}.
\end{proof}

The proof uses nonnegativity of the nodes and their fixed unit weights.
An arbitrary signed or unequal-weight moment problem is not covered by this
argument. In particular, changing the model can change its critical exponent.

\begin{corollary}[Half-length inversion]
\label{cor:half}
For $a,b\in\As$ and $x_i=a_i^2,y_i=b_i^2$,
\begin{equation}
 \norm{a-b}_\infty\le\sqrt8\, E_r(x,y)^{1/(2m+2r)}.
 \label{eq:halfinverse}
\end{equation}
This exponent is also sharp.
\end{corollary}
\begin{proof}
Use $|\sqrt x-\sqrt y|\le\sqrt{|x-y|}$. Sharpness follows from the scaled
positive-node pairs constructed in the next section.
\end{proof}

\section{Exact moment fibres and optimality of the exponents}
\label{sec:fibres}

\subsection{A general construction for every shift}

\begin{lemma}[A moment-preserving tangent]
\label{lem:tangent}
Let $0<u_1<\cdots<u_m<1$ and $P(z)=\prod_i(z-u_i)$. For $r\ge0$ set
\begin{equation}
 b_i=\frac{1}{u_i^rP'(u_i)}.
 \label{eq:tangent}
\end{equation}
The differentials of $p_{r+1},\ldots,p_{r+m-1}$ vanish on $b$, while
$D p_{r+m}(u)b=r+m$. These first $m-1$ differentials are linearly independent.
\end{lemma}
\begin{proof}
Lagrange interpolation of $z^j$, followed by comparison of the coefficient of
$z^{m-1}$, gives
\begin{equation}
 \sum_{i=1}^m\frac{u_i^j}{P'(u_i)}=
 \begin{cases}0,&0\le j<m-1,\\1,&j=m-1.\end{cases}
 \label{eq:lagrange}
\end{equation}
Now $D p_k(u)b=k\sum_i u_i^{k-1-r}/P'(u_i)$, proving the asserted values.
For rank, factor the nonzero $u_i^r$ from each column of the derivative matrix
and the nonzero row constants. What remains is a rectangular Vandermonde
matrix of rank $m-1$. The statements include $m=1$, with no constraints and
an empty rank assertion.
\end{proof}

The implicit function theorem gives a smooth one-dimensional local curve
through $u$, with tangent $b$, on which $p_{r+1},\ldots,p_{r+m-1}$ are exactly
constant. Along this curve $p_{r+m}$ has nonzero derivative. Choose another
point $w$ sufficiently close to $u$ so that all coordinates remain distinct
and in $(0,1)$, and $p_{r+m}(u)\ne p_{r+m}(w)$. Then, for $0<\lambda\le1$,
\begin{align}
 \norm{\lambda u-\lambda w}_\infty
   &=\lambda\norm{u-w}_\infty,\\
 E_r(\lambda u,\lambda w)
   &=\lambda^{m+r}|p_{m+r}(u)-p_{m+r}(w)|.
 \label{eq:sharp-scaled}
\end{align}
A bound with exponent $\alpha>1/(m+r)$ would force a positive constant times
$\lambda$ to be at most a constant times $\lambda^{(m+r)\alpha}$ as
$\lambda\downarrow0$, which is impossible. This proves the sharpness clause
of Theorem~\ref{thm:shifted}. Applying the same pair to half-lengths
$\sqrt\lambda\sqrt{u_i}$ and $\sqrt\lambda\sqrt{w_i}$ proves
Corollary~\ref{cor:half}'s sharpness.

\subsection{An explicit flow that changes only the first visible deficit}
The general construction suffices for sharpness, but statistical variance
recovery needs control of $p_1$ as well. For shift $r=1$, the following
coefficient flow provides it without an unspecified choice of tangent.

\begin{proposition}[First-power-sum flow]
\label{prop:flow}
Let $0<u_1<\cdots<u_m<1$, and let $e_k(0)$ be their elementary symmetric
functions, with $e_0(0)=1$. Define
\begin{equation}
 e_k(s)=\sum_{j=0}^k\frac{s^j}{j!}e_{k-j}(0),\qquad
 P_s(z)=\sum_{k=0}^m(-1)^k e_k(s)z^{m-k}.
 \label{eq:flow}
\end{equation}
For all sufficiently small real $s$, the roots $u_i(s)$ of $P_s$ remain
simple, ordered, and in $(0,1)$. Their power sums satisfy
\begin{align}
 p_1(u(s))&=p_1(u(0))+s,\\
 p_k(u(s))&=p_k(u(0)),\qquad 2\le k\le m,\\
 \frac{\mathrm d}{\mathrm ds}p_{m+1}(u(s))
   &=(-1)^{m+1}(m+1)e_m(s)\ne0.
 \label{eq:flowidentities}
\end{align}
\end{proposition}
\begin{proof}
Simple real roots continue smoothly by the implicit function theorem; choose
$s$ small enough that disjoint root neighborhoods lie in $(0,1)$. Set
$E_s(t)=\sum_{k=0}^m e_k(s)t^k$. Then
\begin{equation}
 E_s(t)=e^{st}E_0(t)\pmod {t^{m+1}},\qquad
 \partial_s E_s(t)=tE_s(t)-e_m(s)t^{m+1}.
\end{equation}
As a formal power series at zero,
\begin{equation}
 \log E_s(t)=\sum_{k\ge1}\frac{(-1)^{k-1}}{k}p_k(u(s))t^k.
\end{equation}
Taking logarithms in the first identity gives the first two assertions.
Dividing the second identity by $E_s(t)$ and taking the coefficient of
$t^{m+1}$ gives the third. Since all roots are positive, $e_m(s)>0$.
\end{proof}

The flow is exact, not just a cancellation of first derivatives. Choosing a
fixed small nonzero $s$ gives two configurations $u,w$ with
\begin{equation}
 p_2(u)=p_2(w),\ldots,p_m(u)=p_m(w),\qquad
 p_1(u)\ne p_1(w),\quad p_{m+1}(u)\ne p_{m+1}(w).
 \label{eq:unknownpair}
\end{equation}
These are the hard pairs used below. They explain concretely how a change in
the unknown Gaussian variance can hide a change in the first power sum while
many additional cumulants also agree.

\section{Identifiability and deterministic Hellinger stability}
\label{sec:hellinger}

For densities $f,g$, use the conventions
\begin{equation}
 H^2(f,g)=\int(\sqrt f-\sqrt g)^2\dd x,\qquad
 \chi^2(f,g)=\int\frac{(f-g)^2}{g}\dd x,
 \quad \TV(f,g)=\frac12\int|f-g|\dd x.
 \label{eq:metrics}
\end{equation}
Here $0\le H^2\le2$, and $\TV\le H$.

\begin{proposition}[Finite-cumulant identifiability]
\label{prop:identifiable}
In the unknown variance model, cumulants of orders
$2,4,6,\ldots,2m+2$ determine $(a,v)$ uniquely. With known variance, orders
$2,4,\ldots,2m$ determine $a$ uniquely.
\end{proposition}
\begin{proof}
Use~\eqref{eq:cumulantidentity} to obtain $p_2,\ldots,p_{m+1}$ when $v$ is
unknown. Theorem~\ref{thm:shifted} with $r=1$ determines the ordered $x$;
then $a_i=\sqrt{x_i}$ and~\eqref{eq:varianceidentity} determines $v$.
For known variance use $r=0$ and $p_1,\ldots,p_m$ instead.
\end{proof}

\begin{theorem}[Global Hellinger inverse moduli]
\label{thm:hellinger}
On the parameter sets of Section~\ref{sec:model},
\begin{align}
 \norm{a-b}_\infty&\le C H(P_{a,v},P_{b,w})^{1/(2m+2)},
 \label{eq:ha}\\
 |v-w|&\le C H(P_{a,v},P_{b,w})^{1/(m+1)}.
 \label{eq:hv}
\end{align}
If $v=w=v_0$ is known, the stronger scale estimate is
\begin{equation}
 \norm{a-b}_\infty\le C H(P_{a,v_0},P_{b,v_0})^{1/(2m)}.
 \label{eq:hknown}
\end{equation}
All three exponents are optimal for uniform bounds on these parameter sets.
\end{theorem}
\begin{proof}[Proof of the upper estimates]
All absolute moments of any fixed order are uniformly bounded in the model.
For raw moments of order $j$, Cauchy--Schwarz gives
\begin{align}
 \abs{\int x^j(f-g)\dd x}
 &\le H(f,g)\left(\int |x|^{2j}(\sqrt f+\sqrt g)^2\dd x\right)^{1/2}\\
 &\le H(f,g)\left(2\E_f|X|^{2j}+2\E_g|X|^{2j}\right)^{1/2}
 \le C_jH(f,g).
 \label{eq:momenthellinger}
\end{align}
Cumulants are polynomials in finitely many raw moments, so they are Lipschitz
on the bounded moment range. Equations~\eqref{eq:varianceidentity}--
\eqref{eq:cumulantidentity} and Theorem~\ref{thm:shifted} now give
\[
 \norm{x-y}_\infty\le C H^{1/(m+1)}
\]
in the unknown case, and the analogous bound with exponent $1/m$ in the
known case. Taking square roots proves the scale estimates. Finally,
\[
 |v-w|\le |\kap_2(P_{a,v})-\kap_2(P_{b,w})|
              +\frac m3\norm{x-y}_\infty
\]
proves~\eqref{eq:hv}, absorbing the term $CH$ since $H\le\sqrt2$.
Optimality follows from the pairs and asymptotics in Sections~\ref{sec:likelihood}
and~\ref{sec:lower}.
\end{proof}

The inverse modulus alone does not prove a minimax sampling lower bound. One
must construct pairs whose \emph{likelihoods}, not merely a chosen finite
list of moments, are close. That step is the subject of the next section.

\section{From exact moment cancellation to likelihood asymptotics}
\label{sec:likelihood}

Write $\phi_v(x)=(2\pi v)^{-1/2}\exp(-x^2/(2v))$. Fix a positive base variance
$v$, let $f_0(x)=\E\phi_v(x-Z)$, and define
\begin{equation}
 J_k(f_0)=\int_\R\frac{(f_0^{(k)}(x))^2}{f_0(x)}\dd x.
 \label{eq:J}
\end{equation}
These quantities are finite despite the unbounded observation space.

\begin{lemma}[Weighted derivative information]
\label{lem:J}
For every integer $k\ge1$,
\begin{equation}
 0<J_k(f_0)\le k!v^{-k}.
 \label{eq:Jbound}
\end{equation}
For $Z=0$ equality holds.
\end{lemma}
\begin{proof}
Differentiating the Gaussian convolution is justified by its integrable
Gaussian bounds. Conditional Cauchy--Schwarz gives pointwise
\[
 \frac{(\E\phi_v^{(k)}(x-Z))^2}{\E\phi_v(x-Z)}
 \le \E\frac{(\phi_v^{(k)}(x-Z))^2}{\phi_v(x-Z)}.
\]
Integrating and using Tonelli bounds $J_k(f_0)$ by the corresponding Gaussian
integral. The derivative identity
\[
 \phi_v^{(k)}(x)=(-1)^k v^{-k/2}\He_k(x/\sqrt v)\phi_v(x)
\]
and the Gaussian orthogonality relation $\E\He_k(G)^2=k!$ evaluate it as
$k!v^{-k}$. Equality for $Z=0$ is immediate. If $J_k(f_0)=0$, its continuous
numerator would vanish identically, so $f_0$ would be a polynomial of degree
at most $k-1$. No nonzero polynomial on $\R$ is an integrable probability
density, a contradiction.
\end{proof}

For any fixed nonnegative vector $u=(u_1,\ldots,u_m)$ put
\begin{equation}
 V_u=\sum_i\sqrt{u_i}U_i,\qquad \beta_u=\frac13p_1(u).
\end{equation}
Two families will be used. The fixed-variance family is
\begin{equation}
 f^{\mathrm{fix}}_{t,u}=\mathrm{density}\bigl(Z+tV_u+\sqrt vG\bigr).
 \label{eq:fixedfamily}
\end{equation}
The compensated family, defined for sufficiently small real $t$, is
\begin{equation}
 f^{\mathrm{comp}}_{t,u}
 =\mathrm{density}\bigl(Z+tV_u+\sqrt{v-t^2\beta_u}\,G\bigr).
 \label{eq:compfamily}
\end{equation}
The total variance contributed by the added variables in the compensated
family is exactly $v$ for every $t$.

\begin{theorem}[Exact leading likelihood term]
\label{thm:likelihood}
Consider either of the following cases:
\begin{enumerate}[label=(\alph*)]
\item The fixed family, with $q\ge1$, and vectors $u,w$ satisfying
$p_k(u)=p_k(w)$ for $1\le k<q$.
\item The compensated family, with $q\ge2$, and vectors $u,w$ satisfying
$p_k(u)=p_k(w)$ for $2\le k<q$.
\end{enumerate}
Assume $d=c_q(p_q(u)-p_q(w))\ne0$, and set $A=d/(2q)!$. Write $f_{t,u},f_{t,w}$
for the chosen family. As $t\to0$,
\begin{equation}
 f_{t,u}-f_{t,w}
   =A t^{2q}f_0^{(2q)}+O_{\mathcal H}(t^{2q+2}),
 \qquad \mathcal H=L^2\bigl(\R,\phi_{3v/4}(x)^{-1}\dd x\bigr).
 \label{eq:densityexpansion}
\end{equation}
Moreover,
\begin{align}
 \chi^2(f_{t,u},f_{t,w})
   &=A^2J_{2q}(f_0)t^{4q}+o(t^{4q}),
 \label{eq:chiexpansion}\\
 H^2(f_{t,u},f_{t,w})
   &=\frac{A^2}{4}J_{2q}(f_0)t^{4q}+o(t^{4q}).
 \label{eq:hexpansion}
\end{align}
In particular, both leading constants are strictly positive.
\end{theorem}

\begin{proof}
We give the tail and denominator details because they are essential for a
statistical lower bound.

\emph{Step 1: a common positive lower envelope.}
Choose $|t|\le t_0$ so that all Gaussian variances in the two families belong
to $[7v/8,v]$. The shifts $W=Z+tV_u$ and $Z+tV_w$ are uniformly bounded by
some $M$. Put $s=3v/4$. For $b\in[7v/8,v]$ and $|W|\le M$, minimization of
a quadratic in $x$ gives
\begin{equation}
 \frac{\phi_b(x-W)}{\phi_s(x)}
 \ge\sqrt{\frac{s}{b}}\exp\left(-\frac{W^2}{2(b-s)}\right)\ge c>0.
 \label{eq:envelope}
\end{equation}
Averaging preserves the inequality. Thus each density in the two families,
and $f_0$, is bounded below by $c\phi_s$ uniformly in $t$.

\emph{Step 2: Taylor expansion in a weighted Hilbert space.}
Every fixed-order derivative in $t$ of
$\phi_{v-t^2\beta_u}(x-Z-tV_u)$, or of its fixed-variance counterpart, is a
finite sum of Gaussian terms times polynomials in $x,Z,V_u$ and bounded
inverse powers of the variance. Consequently, for each derivative order it
is bounded uniformly by
\begin{equation}
 C(1+|x|)^N\exp\left(-\frac{x^2}{2v}+C|x|\right).
 \label{eq:derivativebound}
\end{equation}
After squaring and dividing by $\phi_s$, the quadratic part of the exponent
is $-x^2/(3v)$. This is integrable, including any polynomial and linear
exponential factor. Differentiation under the expectation and differentiation
in $\mathcal H$ are therefore justified to any fixed order. Symmetry of
$V_u$ makes the density an even function of $t$. Taylor's theorem in
$\mathcal H$ gives an expansion with an $O_{\mathcal H}(t^{2q+2})$ remainder
through order $2q$.

\emph{Step 3: identify the first unequal coefficient.}
For each fixed Fourier frequency $\xi$, the logarithm of the uniform part
has its convergent germ
\[
 \log\E e^{i\xi tV_u}
   =\sum_{k\ge1}\frac{c_kp_k(u)}{(2k)!}(i\xi t)^{2k}.
\]
In the compensated family the Gaussian adjustment cancels the $k=1$ term
exactly. The moment equalities in (a) or (b) then show that the first unequal
coefficient of the two characteristic functions is
\[
 A t^{2q}(i\xi)^{2q}\widehat f_0(\xi).
\]
Here the characteristic-function convention is
$\widehat f(\xi)=\int e^{i\xi x}f(x)\dd x$; the Fourier transform of an even
$2q$th derivative has the same multiplier $(i\xi)^{2q}$.
No logarithm of $\widehat f_0$ is needed, so its possible zeros cause no
problem. Coefficients in the Hilbert-space expansion are integrable by
Cauchy--Schwarz against $\phi_s$, and Fourier uniqueness identifies their
difference as $A f_0^{(2q)}$. This proves~\eqref{eq:densityexpansion}.

\emph{Step 4: pass to the likelihood metrics.}
Let $D_t=(f_{t,u}-f_{t,w})/t^{2q}$ and $D=A f_0^{(2q)}$. We know
$D_t\to D$ in $\mathcal H$. The weights $\phi_s/f_{t,w}$ are uniformly bounded
by~\eqref{eq:envelope} and converge pointwise to $\phi_s/f_0$. Weighted
Cauchy--Schwarz first replaces $D_t$ by $D$ in the integral, and dominated
convergence then gives
\[
 \int\frac{D_t^2}{f_{t,w}}\dd x\longrightarrow\int\frac{D^2}{f_0}\dd x.
\]
This is~\eqref{eq:chiexpansion}. Finally, the identity
\[
 (\sqrt f-\sqrt g)^2=\frac{(f-g)^2}{(\sqrt f+\sqrt g)^2}
\]
has denominator converging pointwise to $4f_0$ and bounded below by
$4c\phi_s$. The identical argument gives~\eqref{eq:hexpansion}. Positivity
of the limits follows from Lemma~\ref{lem:J}.
\end{proof}

This proof explains why merely matching finitely many moments is not enough:
the required control is in a space weighted by a reciprocal density. The
positive Gaussian component supplies a common tail envelope and prevents
rare tail events from defeating the lower-bound construction.

\section{A feasible moment estimator and upper bounds}
\label{sec:upper}

This section constructs a measurable estimator and proves the upper half of
Theorem~\ref{thm:main}. The procedure is deliberately elementary and globally
specified. It is not claimed to be computationally optimal.

\subsection{Uniform moment control and cumulant estimation}
For every integer $j\ge1$, all parameters in the unknown model satisfy
\begin{equation}
 \E|X_{a,v}|^j\le B_j
 :=2^{j-1}\left((L+m)^j+\overline v^{j/2}\E|G|^j\right).
 \label{eq:Bj}
\end{equation}
For known variance use $v_0$ instead of $\overline v$. Given observations
$X_1,\ldots,X_n$, set $\widehat\mu_0=1$ and
\begin{equation}
 \widehat\mu_j=
 \clip_{[-B_j,B_j]}\left(\frac1n\sum_{\ell=1}^n X_\ell^j\right).
 \label{eq:rawestimator}
\end{equation}
Because the true raw moment lies in the clipping interval,
\begin{equation}
 \E|\widehat\mu_j-\mu_j|^2\le\frac{B_{2j}}{n}.
 \label{eq:rawMSE}
\end{equation}
Cumulants are computed recursively from these clipped moments:
\begin{equation}
 \widehat\kap_j=\widehat\mu_j-
 \sum_{\ell=1}^{j-1}\binom{j-1}{\ell-1}
                \widehat\kap_\ell\widehat\mu_{j-\ell}.
 \label{eq:cumulantrecursion}
\end{equation}
The exact identity follows by differentiating the relation between a moment
generating function and its logarithm, then comparing coefficients. Its
estimated version is a polynomial in a bounded moment vector. Every fixed
polynomial is Lipschitz on this compact box, so for fixed $K$,
\begin{equation}
 \E\max_{1\le j\le K}|\widehat\kap_j-\kap_j(X_{a,v})|^2\le\frac{C_K}{n}.
 \label{eq:cumulantMSE}
\end{equation}
The estimators need not be unbiased; the bound, not unbiasedness, is what the
inversion argument requires.

For unknown variance use $K=2m+2$ and define
\begin{equation}
 \widehat p_k=\frac{\widehat\kap_{2k}-\kap_{2k}(Z)}{c_k},
 \qquad 2\le k\le m+1.
 \label{eq:phatunknown}
\end{equation}
For known variance use $K=2m$, the same formula for $2\le k\le m$, and
\begin{equation}
 \widehat p_1=3\bigl(\widehat\kap_2-\kap_2(Z)-v_0\bigr).
 \label{eq:phatknown}
\end{equation}
With $r=1$ in the unknown case and $r=0$ in the known case, write
\begin{equation}
 \Delta=\max_{r+1\le k\le r+m}|\widehat p_k-p_k(x)|.
\end{equation}
Equations~\eqref{eq:cumulantMSE} and $c_k\ne0$ imply
\begin{equation}
 \E\Delta^2\le C/n,\qquad \E\Delta\le Cn^{-1/2}.
 \label{eq:delta}
\end{equation}

\subsection{Globally feasible fitting instead of unstable root clipping}
A noisy moment vector need not correspond to nonnegative real nodes. We
therefore fit within the admissible set rather than form a polynomial and
silently discard its complex or negative roots. Let
\begin{equation}
 \mathcal G_n=\Xs\cap\{0,1/n,\ldots,1\}^m.
\end{equation}
Choose the lexicographically first minimizer
\begin{equation}
 \widehat x\in\operatorname*{arg\,min}_{z\in\mathcal G_n}
 \max_{r+1\le k\le r+m}|p_k(z)-\widehat p_k|,
 \qquad \widehat a_i=\sqrt{\widehat x_i}.
 \label{eq:fit}
\end{equation}
There are exactly $\binom{n+m}{m}$ grid points, so the rule exists and is
measurable. It is an exhaustive finite algorithm when the required constants
and objective comparisons are provided to the desired accuracy. Its grid
size is a statistical existence device, not a recommendation for large-$m$
numerical implementation.

Coordinatewise downward rounding of the true $x$ to this grid preserves
order and changes each coordinate by at most $1/n$. Since $z\mapsto z^k$ is
$k$-Lipschitz on $[0,1]$, the rounded vector has objective at most
$\Delta+m(m+r)/n$. The minimizer therefore satisfies
\begin{equation}
 E_r(\widehat x,x)\le 2\Delta+\frac{m(m+r)}n.
 \label{eq:fitresidual}
\end{equation}
Theorem~\ref{thm:shifted} and Corollary~\ref{cor:half} yield
\begin{align}
 \norm{\widehat x-x}_\infty
 &\le C(\Delta+n^{-1})^{1/(m+r)},
 \label{eq:xestbound}\\
 \norm{\widehat a-a}_\infty
 &\le C(\Delta+n^{-1})^{1/(2m+2r)}.
 \label{eq:aestbound}
\end{align}
Jensen's inequality and~\eqref{eq:delta} prove
\begin{equation}
 \sup\E\norm{\widehat a-a}_\infty\le Cn^{-1/(4m+4r)}.
 \label{eq:aupper}
\end{equation}
Taking $r=0,1$ gives the two scale upper bounds.

In the unknown case define the variance estimator
\begin{equation}
 \widehat v=\clip_{[\underline v,\overline v]}
 \left(\widehat\kap_2-\kap_2(Z)-\frac13p_1(\widehat x)\right).
 \label{eq:vestimator}
\end{equation}
Clipping cannot increase its distance from the true $v$. Thus
\begin{equation}
 \E|\widehat v-v|
 \le\E|\widehat\kap_2-\kap_2(X)|
        +\frac m3\E\norm{\widehat x-x}_\infty
 \le C n^{-1/(2m+2)},
 \label{eq:vupper}
\end{equation}
where the faster $n^{-1/2}$ term is absorbed. This completes all upper bounds
in Theorem~\ref{thm:main}.

\subsection{Approximate optimization and confidence bounds}
The estimator remains valid if the achieved objective is within $\eta_n$ of
the global grid minimum. Equation~\eqref{eq:fitresidual} then acquires an
additional $\eta_n$. In particular, a certified objective gap
$\eta_n=O(n^{-1/2})$ preserves the rates. An arbitrary local optimum, without
a bound on its gap, has no such guarantee.

For $0<\delta<1$, Markov's inequality applied to $\Delta^2$ gives, with
probability at least $1-\delta$,
\begin{equation}
 \norm{\widehat a-a}_\infty
 \le C\left((n\delta)^{-1/2}+n^{-1}\right)^{1/(2m+2r)}.
 \label{eq:confidence}
\end{equation}
These elementary confidence bounds are conservative. Sharper concentration
or adaptive confidence procedures are separate questions; neither is needed
for the minimax mean-risk result.

\section{Matching statistical lower bounds}
\label{sec:lower}

We first record the testing reduction with the metric conventions used here.
If two parameters are separated by $D$ in a metric loss, the test that chooses
the parameter nearer an estimator has an error only if that estimator's loss
is at least $D/2$. Consequently, for every estimator,
\begin{equation}
 \max_{i=0,1}\E_i[\mathrm{loss}]
 \ge\frac D4\left(1-\TV(P_0^{\otimes n},P_1^{\otimes n})\right).
 \label{eq:lecam}
\end{equation}
This follows because the minimum sum of the two simple testing errors is
$1-\TV$. Hellinger affinity tensorizes, so
\begin{equation}
 H^2(P^{\otimes n},Q^{\otimes n})
 =2\left[1-\left(1-\frac{H^2(P,Q)}2\right)^n\right]
 \le nH^2(P,Q).
 \label{eq:tensor}
\end{equation}
Together with $\TV\le H$, this controls the bracket in~\eqref{eq:lecam}.

\subsection{Unknown Gaussian variance}
Choose $v_*\in(\underline v,\overline v)$ and the fixed positive distinct
vectors $u,w$ from Proposition~\ref{prop:flow}, with a sufficiently small
nonzero flow parameter. They satisfy~\eqref{eq:unknownpair}. For small $t>0$
define two model parameters by
\begin{align}
 a_i(t)&=t\sqrt{u_i},& v_u(t)&=v_* -\frac{t^2}{3}p_1(u),\\
 b_i(t)&=t\sqrt{w_i},& v_w(t)&=v_* -\frac{t^2}{3}p_1(w).
 \label{eq:hardunknown}
\end{align}
They belong to the parameter set for every sufficiently small $t$, including
the variance interval. Their scale and variance separations are
\begin{equation}
 \norm{a(t)-b(t)}_\infty=D_a t,\qquad
 |v_u(t)-v_w(t)|=D_v t^2,
 \label{eq:separations}
\end{equation}
where $D_a,D_v>0$ are fixed. Theorem~\ref{thm:likelihood}, with $q=m+1$,
gives
\begin{equation}
 H^2(P_{a(t),v_u(t)},P_{b(t),v_w(t)})
 =K t^{4m+4}+o(t^{4m+4}),\qquad K>0.
 \label{eq:hardHunknown}
\end{equation}
Set $t=\gamma n^{-1/(4m+4)}$ with fixed $\gamma>0$ small enough. Equations
\eqref{eq:tensor} and~\eqref{eq:hardHunknown} ensure product total variation
at most $1/2$ for all sufficiently large $n$. Applying~\eqref{eq:lecam}
separately to the two separations in~\eqref{eq:separations} proves
\begin{equation}
 \risk^a_n\ge c n^{-1/(4m+4)},\qquad
 \risk^v_n\ge c n^{-1/(2m+2)}.
\end{equation}
The same pair proves both lower bounds; the variance difficulty is not an
artifact of the plug-in estimator.

\subsection{Known Gaussian variance}
Use the construction of Lemma~\ref{lem:tangent} with $r=0$ to choose positive
vectors $u,w$ for which $p_1,\ldots,p_{m-1}$ agree but $p_m$ differs. For $m=1$
there are no matching constraints, and any two distinct positive nodes
suffice. Set $a=t\sqrt u$, $b=t\sqrt w$, and keep the Gaussian variance
fixed at $v_0$. Theorem~\ref{thm:likelihood}, now with $q=m$ in the fixed
family, gives $H^2\asymp t^{4m}$. Choosing $t=\gamma n^{-1/(4m)}$ in the same
two-point argument proves
\begin{equation}
 \risk^{a,\mathrm{known}}_n\ge c n^{-1/(4m)}.
\end{equation}
This completes Theorem~\ref{thm:main}.

\subsection{Optimality of the deterministic moduli}
The same pairs also complete Theorem~\ref{thm:hellinger}. In the unknown
case $H\asymp t^{2m+2}$ while scale distance is $D_at$ and variance distance
is $D_vt^2$. An exponent larger than $1/(2m+2)$ or $1/(m+1)$, respectively,
would contradict these asymptotics. In the known case $H\asymp t^{2m}$ and
scale distance is proportional to $t$, excluding a larger exponent there.

\subsection{Why these are global, not everywhere-local, rates}
The hard parameters approach the boundary where all uniform factors vanish.
They do not say that well-separated positive factors are always recovered
slowly. In fact, for the unknown-variance moment map
$F(x)=(p_2(x),\ldots,p_{m+1}(x))$,
\begin{equation}
 \det DF(x)=(m+1)!\left(\prod_{i=1}^m x_i\right)
                   \prod_{1\le i<j\le m}(x_j-x_i).
 \label{eq:jacobian}
\end{equation}
To check this, factor $k$ from the row for $p_k$ and $x_i$ from column $i$;
the remaining determinant is Vandermonde. At a strictly positive, distinct
configuration this determinant is nonzero. The inverse function theorem,
together with the moment bounds of Section~\ref{sec:upper}, yields local
$n^{-1/2}$ estimation on a sufficiently small regular neighborhood, including
for $v$ and for the half-lengths. The slow global exponents reflect loss of
rank near vanishing or colliding configurations, not a uniform penalty at
every point.

\section{Detection, rank, and the distinction from recovery}
\label{sec:testing}

We prove Theorem~\ref{thm:detectionintro}. For unknown variance the exact
fourth-cumulant contrast is
\begin{equation}
 T=\kap_4(X_{a,v})-\kap_4(Z)
   =-\frac{2}{15}\sum_i a_i^4.
 \label{eq:testcontrast}
\end{equation}
Let $\widehat T=\widehat\kap_4-\kap_4(Z)$, using the clipped-moment procedure.
By~\eqref{eq:cumulantMSE}, $\E|\widehat T-T|^2\le C/n$ uniformly. Reject
$a=0$ when
\begin{equation}
 \widehat T<-h^4/15.
\end{equation}
Under the null $T=0$; under any alternative with $\norm a_\infty\ge h$,
$T\le-2h^4/15$. Either kind of error requires
$|\widehat T-T|\ge h^4/15$. Hence the sum of maximal errors is at most
$C/(nh^8)$.

For the converse, fix $v_*\in(\underline v,\overline v)$ and compare
\begin{equation}
 (a,v)=(0,v_*)\quad\text{with}\quad
 (a,v)=((0,\ldots,0,h),v_*-h^2/3).
 \label{eq:detectionpair}
\end{equation}
For small $h$ both are admissible. Apply Theorem~\ref{thm:likelihood} with
$q=2$, one node equal to $1$ versus all nodes zero, and $t=h$. There are no
higher matching constraints before $p_2$; compensation matches variance.
The one-observation squared Hellinger distance is $O(h^8)$. If $nh_n^8\to0$,
product total variation tends to zero by~\eqref{eq:tensor}, so every test
has the sum of its two errors at least $1-o(1)$. Composite maximal risks
are at least these two simple risks.

When variance is known, use
\begin{equation}
 T_2=\kap_2(X_{a,v_0})-\kap_2(Z)-v_0
     =\frac13\sum_i a_i^2
\end{equation}
and reject for $\widehat T_2>h^2/6$. This gives $C/(nh^4)$ by the same mean
square argument. The fixed-variance one-factor pair has $H^2=O(h^4)$ by
Theorem~\ref{thm:likelihood} with $q=1$, proving the lower bound.

At a fixed large multiple of the critical scale the upper bounds can be
made arbitrarily small by increasing that multiple. Consistency, meaning
errors tending to zero with $n$, follows for diverging multiples. We do not
confuse these two statements.

\begin{corollary}[No uniform exact factor-count estimation without separation]
\label{cor:rank}
Let the factor count be $\#\{i:a_i>0\}$. Over either full parameter set of
Section~\ref{sec:model}, no sequence of estimators of this count has
probability of error tending to zero uniformly over all parameters.
\end{corollary}
\begin{proof}
Choose $h_n$ below the respective detection scale in~\eqref{eq:detectionpair},
or in its known-variance counterpart. The two parameters have different
counts, but their product laws have total variation tending to zero. An
exact count estimator would distinguish them, contradicting the testing
lower bound.
\end{proof}

This does not rule out pointwise consistency when every nonzero factor has
a fixed positive half-length. It rules out \emph{uniform} recovery of the
count without a lower bound on the smallest positive half-length.

\section{Explicit two-factor and one-factor examples}
\label{sec:examples}

\subsection{A two-factor pair with a sixth-order first difference}
Take squared-scale templates
\begin{equation}
 u=(0,5),\qquad w=(3,4).
 \label{eq:explicitpair}
\end{equation}
The templates themselves exceed $1$, but the actual squared scales are
$t^2u,t^2w$, which lie in $[0,1]^2$ for sufficiently small $t$. They satisfy
\begin{equation}
 p_1(u)=5,\quad p_1(w)=7,\qquad
 p_2(u)=p_2(w)=25,\qquad p_3(u)-p_3(w)=34.
\end{equation}
Choose Gaussian variances $v_*-5t^2/3$ and $v_*-7t^2/3$. The two uniform
half-length vectors are $(0,\sqrt5\,t)$ and $(\sqrt3\,t,2t)$, separated by
$\sqrt3\,t$, while the variances differ by $2t^2/3$. The first unequal added
cumulant is the sixth:
\begin{equation}
 \Delta\kap_6=\frac{16}{63}\,34t^6=\frac{544}{63}t^6.
\end{equation}
Theorem~\ref{thm:likelihood} therefore gives
\begin{equation}
 f_{t,u}-f_{t,w}=\frac{34}{2835}t^6 f_0^{(6)}+O_{\mathcal H}(t^8).
 \label{eq:explicitdensity}
\end{equation}
For $Z=0$ and $v_*=1$, the corresponding chi-square expansion is
\begin{equation}
 \chi^2(f_{t,u},f_{t,w})
 =\frac{18496}{178605}t^{12}+o(t^{12}).
 \label{eq:explicitchi}
\end{equation}
Indeed, $J_6(\phi_1)=6!=720$, and
$720(34/2835)^2=18496/178605$. This gives the unknown-variance two-factor
scale $t\asymp n^{-1/12}$ and variance separation of order $n^{-1/6}$ without
requiring a numerically solved polynomial.

For a known-variance comparison, use $u=(1,4)$ and $w=(2,3)$.
Their first power sums agree and their second power sums differ by $4$.
The leading density difference is
\begin{equation}
 f_{t,u}^{\mathrm{fix}}-f_{t,w}^{\mathrm{fix}}
   =-\frac1{45}t^4f_0^{(4)}+O_{\mathcal H}(t^6),
\end{equation}
which gives the faster two-factor scale $n^{-1/8}$.

\subsection{One small factor and an unknown Gaussian variance}
Compare $\sqrt{v_*}G$ with $tU+\sqrt{v_*-t^2/3}G$. Variance is exactly
matched, and the fourth cumulant is the first non-Gaussian signal:
\begin{equation}
 f_t-\phi_{v_*}=-\frac{t^4}{180}\phi_{v_*}^{(4)}+O_{\mathcal H}(t^6),
 \qquad
 \chi^2(f_t,\phi_{v_*})
   =\frac{4!}{180^2v_*^4}t^8+o(t^8).
\end{equation}
This is the simplest instance of the $n^{-1/8}$ threshold. If the variance
cannot adjust, the leading density difference is instead
$t^2\phi_{v_*}''/6$, giving an $n^{-1/4}$ threshold. The nuisance parameter
removes a genuine leading observable term.

\subsection{A Fourier diagnostic with a known exact limit}
For the pair~\eqref{eq:explicitpair}, $Z=0$, $v_*=1$, and frequency $1$, put
\begin{align}
 \ell_u(t)&=-\frac{1-5t^2/3}{2}+\log\sinc(\sqrt5\,t),\\
 \ell_w(t)&=-\frac{1-7t^2/3}{2}
                 +\log\sinc(\sqrt3\,t)+\log\sinc(2t).
\end{align}
For small positive $t$ the sinc factors are positive, so these real logarithms
are unambiguous. The proved coefficient identity gives
\begin{equation}
 \frac{\ell_u(t)-\ell_w(t)}{t^6}\longrightarrow-\frac{34}{2835}.
\end{equation}
The included program evaluates this diagnostic at 80-digit working precision:
\begin{center}
\begin{tabular}{@{}rr@{}}
\toprule
$t$ & $(\ell_u(t)-\ell_w(t))/t^6$\\
\midrule
$0.2$ & $-0.0123041895267270$\\
$0.1$ & $-0.0120695349593127$\\
$0.05$ & $-0.0120120178007690$\\
$0.025$ & $-0.0119977087830769$\\
$0.0125$ & $-0.0119941358994491$\\
\bottomrule
\end{tabular}
\end{center}
The calculation is a diagnostic, not an interval-certified proof of the
asymptotic statement or of the likelihood integrals.

\section{Verification record and a formalization route}
\label{sec:verification}

\subsection{What the accompanying program checks}
The package contains \texttt{verify\_results.py}, requiring SymPy and mpmath,
with random seed \texttt{20260929}. It produces a machine-readable JSON record
and a plain-text summary. The completed checks are:
\begin{enumerate}
\item $7000$ exact rational instances of~\eqref{eq:shifted}, for
$1\le m\le7$, $0\le r\le3$, including zero entries and repeated nodes.
Each comparison uses integer or rational arithmetic, not floating-point
H\"older powers.
\item $40$ exact Lagrange/Vandermonde tangent checks for $1\le m\le8$ and
$0\le r\le4$.
\item Eight exact coefficient-flow checks, one for each $1\le m\le8$.
For every case, a nonzero rational shift is accompanied by $m$ disjoint
positive rational intervals on each of which the resulting degree-$m$
polynomial changes sign. This certifies all its roots to be real, positive,
and distinct. Power-sum conservation and the derivative of $p_{m+1}$ are
verified symbolically.
\item Exact evaluation of the two-factor leading constants in
\eqref{eq:explicitdensity}--\eqref{eq:explicitchi}, followed by the
high-precision Fourier diagnostic above.
\end{enumerate}
Every listed exact check passed. These finite checks support the algebra and
examples but do not replace the all-parameter proofs. The program does not
claim a simulated minimax theorem, an empirical risk study, an interval bound
on an integral, or a proof-assistant verification.

\subsection{A useful division of the formalization work}
The argument separates into three fairly independent layers.

\emph{Finite algebra and order.} Formalize ordered nonnegative vectors,
power sums, counting functions, their sign blocks, and
Theorem~\ref{thm:shifted}. The bottleneck is not symbolic manipulation but the
finite combinatorial bound on sign changes, including ties. Lagrange
interpolation and the elementary-symmetric flow can then be proved as exact
polynomial identities. The program's rational root intervals are natural
certificate data for this layer, provided a proved checker supplies their
semantics.

\emph{Analysis of Gaussian convolutions.} Formalize the common lower envelope,
Gaussian derivative/Hermite bounds, parameter differentiation in a weighted
$L^2$ space, and the passage from that convergence to chi-square and Hellinger
limits. This is the essential analytic bridge; an unproved assertion that
moment matching implies Hellinger closeness must not substitute for it.

\emph{Statistical decisions.} Formalize measurable finite-grid minimization,
uniform moment bounds, clipping, the two-point testing inequality, and
Hellinger tensorization. Once these are available, the exponents follow by
short substitutions. Theorems about the known and unknown variance experiments
should retain different parameter types or explicit hypotheses, preventing
an accidental removal of the positive-smoothing assumption.

This is a roadmap, not delivered Lean code. The repository's existing formal
Fabius tools motivate the background, but no assertion in this article is
promoted to a formal theorem merely because an adjacent repository module
has been kernel-checked.

\section{Further research questions and success criteria}
\label{sec:future}

\begin{problem}[Remove Gaussian smoothing without losing the lower bound]
For the unsmoothed model $Z_{\mathrm{up}}+\sum_i a_iU_i$, determine the sharp
sampling minimax rates at vanishing factors. The key issue is whether a
suitable replacement for the weighted Taylor argument holds near the flat,
moving support boundary. A first concrete target is finiteness, or a precise
failure mode, of $\int (g^{(k)})^2/g$ for the up density $g$, together with
control of the support change under convolution. A total-variation inverse
modulus alone does not settle this problem.
\end{problem}

\begin{problem}[Classify all local strata with unknown variance]
At a base configuration with positive squared-scale multiplicities
$m_1,\ldots,m_s$ and $m_0$ zero slots, determine the exact local pairwise
Hellinger inverse exponent. A candidate denominator for the half-length
exponent is
\begin{equation}
 M=\max\bigl(\{m_1,\ldots,m_s\}\cup
                  \{2(m_0+1):m_0>0\}\bigr).
 \label{eq:localconjecture}
\end{equation}
The proposed exponent is $1/M$ for the Hellinger modulus and $1/(2M)$ for
the corresponding sampling rate. This is a conjectural extension, not a
result proved here. It agrees with the fully vanishing and regular cases
proved above, but mixed clusters require a separate local argument and
matching likelihood constructions.
\end{problem}

\begin{problem}[Adapt to collision and vanishing patterns]
Construct a single practical estimator attaining the regular $n^{-1/2}$ rate
away from singularities and the correct slower rates on every singular
stratum without being told the stratum. Determine which honest confidence
sets can adapt to these rates. Separate impossibility due to testing between
strata from failures of a particular algorithm.
\end{problem}

\begin{problem}[Efficient globally certified shifted-moment fitting]
Replace the $\binom{n+m}{m}$ grid search by an algorithm with explicit
complexity in $m$ and the requested objective tolerance. One route is to use
hyperbolic polynomial coefficients and the special flow of
Proposition~\ref{prop:flow}; another is constrained optimization with a
certificate of a global objective gap. A useful endpoint is a solver whose
certificate guarantees the $O(n^{-1/2})$ gap sufficient for the statistical
rates, including multiple and zero roots.
\end{problem}

\begin{problem}[Growing capacity and the infinite-factor transition]
Allow $m=m_n\to\infty$ or a summable infinite half-length spectrum. Track
constants in the shifted-moment, cumulant-estimation, and likelihood bounds
instead of hiding them in fixed-$m$ notation. Determine where the polynomial
rates proved here cease to be informative and connect that regime to the
repository's infinite-factor instability questions. The fixed-$m$ theorem
cannot simply be read with a growing $m$ substituted into its exponent.
\end{problem}

\begin{problem}[Learn the background from calibration samples]
Suppose the background law is not known exactly but a second independent
sample from $Z$ is available. Quantify the joint dependence on the observation
sample size and calibration sample size, particularly when high cumulants of
$Z$ must be estimated. A second variant is an explicitly bounded deterministic
error in each background cumulant. Determine when background uncertainty
becomes the dominant obstruction.
\end{problem}

\begin{problem}[Multivariate uniform directions and covariance confounding]
Replace each $a_iU_i$ by $U_i b_i$ with unknown vectors $b_i\in\R^d$, and let
the Gaussian covariance matrix be unknown. Higher even cumulants are then
sums of tensors $b_i^{\otimes 2k}$. Establish identifiability up to signs and
permutations, and sharp rates near collinear, colliding, or vanishing vectors.
The one-dimensional sign-count proof no longer applies directly; it isolates
what a genuine multivariate replacement must accomplish.
\end{problem}

\begin{problem}[Optimal constants and different moment windows]
Determine the best constant in~\eqref{eq:shifted}, including its dependence
on $m,r$. More structurally, classify which nonconsecutive power windows
retain injectivity and what their optimal inverse exponents are. Then examine
known unequal positive weights in the power sums. Equal unit weights were
used in the excursion-count argument, so a weighted extension needs its own
hypotheses rather than an automatic transfer.
\end{problem}

\section{Conclusion}
Unknown Gaussian variance does more than add an extra parameter. It absorbs
the second-cumulant signature of small uniform factors, so the recoverable
power-sum window shifts from $p_1,\ldots,p_m$ to $p_2,\ldots,p_{m+1}$.
The sharp shifted-moment inequality quantifies that loss algebraically;
exact moment fibres and a weighted Gaussian likelihood expansion show that
it persists at the level of statistical information. The resulting scale
rates are $n^{-1/(4m)}$ with known variance and $n^{-1/(4m+4)}$ with unknown
variance, with variance estimation at $n^{-1/(2m+2)}$. Detection uses only the
first unconfounded non-Gaussian cumulant and has the different, capacity-free
threshold $n^{-1/8}$.

These conclusions extend a specific ProveIt research direction without
assuming the unreviewed preceding report as a proved analytic black box.
They leave explicit boundaries: positive smoothing, fixed capacity, known
bounded background, and conventional rather than machine-checked proofs.
The most direct next advance is a sharp treatment of mixed local strata or
of the unsmoothed up-law likelihood boundary.

\begin{thebibliography}{9}
\bibitem{Arias}
J.~Arias de Reyna,
\emph{An infinitely differentiable function with compact support: Definition
and properties}, arXiv:1702.05442 (2017).
English translation of the 1982 article in
\emph{Rev. Real Acad. Ciencias Madrid} \textbf{76}, 21--38.
\url{https://arxiv.org/abs/1702.05442}.

\bibitem{BatenkovYomdin}
D.~Batenkov and Y.~Yomdin,
\emph{Geometry and Singularities of the Prony mapping},
arXiv:1301.1336 (2013).
\url{https://arxiv.org/abs/1301.1336}.

\bibitem{HeinrichKahn}
P.~Heinrich and J.~Kahn,
\emph{Optimal rates for finite mixture estimation},
arXiv:1507.04313 (2015).
\url{https://arxiv.org/abs/1507.04313}.

\bibitem{WuYang}
Y.~Wu and P.~Yang,
\emph{Optimal estimation of Gaussian mixtures via denoised method of moments},
arXiv:1807.07237v2 (2019).
\url{https://arxiv.org/abs/1807.07237v2}.

\bibitem{RepoFinite}
V.~Reshetnikov, ProveIt repository research collection,
\emph{Sharp Stability Strata for Finite Fabius--Rvachev Deconvolution},
article dated 28 September 2026, indexed as an unreviewed intake on
29 September 2026. Read at commit
\texttt{c130dba623c420551d90db41dae7b3d9cff72dfd}.
Repository \url{https://github.com/VladimirReshetnikov/ProveIt}; path
\path{Analysis/FabiusFunction/docs/semi-formalized-research-frontiers/drafts/inverse-and-sampling/Sharp_Stability_Strata_Fabius_Rvachev_Deconvolution/article.tex}.

\bibitem{RepoIndex}
V.~Reshetnikov, ProveIt repository,
\emph{Inverse and sampling: research index}, read 29 September 2026 at the
same commit. Repository \url{https://github.com/VladimirReshetnikov/ProveIt};
path
\path{Analysis/FabiusFunction/docs/semi-formalized-research-frontiers/drafts/inverse-and-sampling/README.md}.
\end{thebibliography}
\end{document}
